\documentclass[11pt,aspectratio=169]{beamer}
% Change text margins
\setbeamersize{text margin left = 12pt, text margin right = 12pt} 
\setbeamercovered{transparent}

\settowidth{\leftmargini}{\usebeamertemplate{itemize item}}
\addtolength{\leftmargini}{\labelsep}

% \begin{columns} -----------------------------------
\usepackage{multicol}
\usepackage{natbib}
\usepackage{amsmath}
\usepackage{amsfonts}
\usepackage{verbatim}
\usepackage{mathrsfs}
\usepackage{amssymb}
\usepackage{amsthm}
\usepackage{bbm}
\usepackage{bm}
\usepackage{booktabs}
\usepackage{longtable}
\usepackage{multirow}
\usepackage{bigdelim}
%\onehalfspacing
\usepackage{courier}

% Title --------------------------------------------
\title{A General Equilibrium Analysis of Rent Control with Heterogeneous Households}
\date{\today}
\author{Shahar Rotberg$^*$ and Lin Zhang$\dagger$ (CMHC)}

\begin{document}

% ------------------------------------------------------------------------------
\begin{frame}
\maketitle
 
 \vspace{1cm}
{\footnotesize $^*$Shahar Rotberg co-authored this paper in his own personal capacity and time. $\dagger$The opinions expressed herein do not reflect the views of the CMHC.}
\end{frame}

%%%%%%%%%%%%%%%%%%%%%%%%%
%%%%%%%%%%%%%%%%%%%%%%%%%
\begin{frame}{Introduction}{Research questions}

How does rent control affect eq'm rents, house prices, and welfare in general equilibrium?

\vspace{0.5cm}

Who losses and who benefits from rent control?

\vspace{0.5cm}

How do transitions out of rent control look like?
 
\end{frame}

\begin{frame}{Introduction}{Motivation}

\uncover<1-4>{In this work, we argue that:}
\begin{enumerate}
    \item \uncover<2>{\textcolor{blue}{Although in long-run rent control reduces avg. rent it reduces welfare}}
    \item \uncover<3>{\textcolor{blue}{Fully repealing rent control raises welfare in long-run and immediately upon implementation}}
    \item \uncover<4>{\textcolor{blue}{No amount of rent control is welfare-maximizing in long-run}}
\end{enumerate}
    
\end{frame}

%%%%%%%%%%%%%%%%%%%%%%%%%
%%%%%%%%%%%%%%%%%%%%%%%%%
\begin{frame}{Introduction}{Motivation}

Many cities have considered or have implemented rent controls \citep{fraser-1981}

\vspace{0.5cm}

Last 30 years focused on rent controls with tenancy deregulation and exemptions
\begin{itemize}
    \item Newly-built units are exempted from rent control
    \item Controlled units must follow rent increase guidelines 
    \begin{itemize}
        \item Landlords can reset rental prices with new tenants
    \end{itemize}
\end{itemize}

\vspace{0.5cm}

Scant general eq'm analysis of rent control \citep{diamond-2019}
\begin{itemize}
    \item \citet{hayek-1929} and \citet{fs-1946} study rent freeze's effect on supply
    \item  Economics textbooks analyze its dead-weight loss
    \item \citet{fallis-1985} study controls with exemptions using partial eq'm model
    \item \citet{fmn-2022} is related but more similar to social affordable housing
\end{itemize}
\end{frame}

{\setbeamercolor{background canvas}{bg=black}
\begin{frame}{}

\centering \Large
\textcolor{white}{What we do}
    
\end{frame}

}

%%%%%%%%%%%%%%%%%%%%%%%%%
%%%%%%%%%%%%%%%%%%%%%%%%%
\begin{frame}{Introduction}{What we do (part 1)}

Model rent control with exemptions in general eq'm with heterogeneous hhlds

\begin{itemize}
    \item Households choose to own or rent
    \item Rental market has controlled and uncontrolled segments
    \begin{itemize}
        \item Real rents in controlled units decline the longer tenants remain in same unit  
    \end{itemize}
    \item Renters optimally choose how long to stay in same unit
    \item Renters optimally choose in which market to rent
    \begin{itemize}
        \item To achieve eq'm, renters in controlled market pay initially more than in free market
        \item Renters planning to stay longer in same unit sort to controlled market
        \begin{itemize}
            \item High-income renters plan to buy early on $\rightarrow$ sort to free market
            \item Low-income renters plan to remain renters longer $\rightarrow$ sort to controlled market
        \end{itemize}
    \end{itemize}
\end{itemize}

\end{frame}

\begin{frame}{Introduction}{What we do (part 2)}

Calibrate model to data from Toronto where markets are partly controlled and partly free:

\begin{itemize}
    \item Census, 2001-2016
    \item Rental Market Survey, 1992-2021
    \item Survey of Financial Security, 2016
\end{itemize}

\vspace{0.5cm}

Calibrated model imitates reasonably well data from Toronto on:

\begin{itemize}
    \item Life-cycle homeownership and debt
    \item Sorting behavior into free and controlled markets
    \item Wealth inequality
    \item Other aggregate statistics
\end{itemize}
    
\end{frame}


\begin{frame}{Introduction}{What we do (part 3)}

We conduct 4 experiments and compare to our benchmark:

\begin{itemize}
    \item Long-run full decontrol
    \item Transition from benchmark to long-run full decontrol - modelled as permanent “MIT shock”
    \item Long-run Full control
    \item Optimal degree of rent control in long-run
    \begin{itemize}
        \item Maximize ex-ante expected utility of newborns \citep{conesa-et-al,gkkoc}
    \end{itemize}
\end{itemize}
    
\end{frame}

{\setbeamercolor{background canvas}{bg=black}
\begin{frame}{}

\centering \Large
\textcolor{white}{Main results}
    
\end{frame}

}

%%%%%%%%%%%%%%%%%%%%%%%%%
%%%%%%%%%%%%%%%%%%%%%%%%%
\begin{frame}{Introduction}{Main results (part 1)}

In long-run, full rent decontrol would:

\begin{itemize}
    \item Raise average rents but reduce free market rents and rents of new tenants
    \begin{itemize}
        \item Decontrol $\rightarrow$ less benefit to renting longer $\rightarrow$ more homeownership $\rightarrow$ less rental demand
        \item Lower rents of new tenants would benefit young, low-income households 
    \end{itemize}
\end{itemize}

\vspace{0.5cm}

In short-run, full decontrol would:

\begin{itemize}
    \item Increase welfare upon implementation
    \begin{itemize}
        \item Renters in large controlled units switch to smaller units $\rightarrow$ reduce rental demand immediately
        \item Lower free market rents $\rightarrow$ lower initial rents of new tenants upon implementation
    \end{itemize}
    \item But raise average rents
\end{itemize}

\end{frame}

\begin{frame}{Introduction}{Main results (part 2)}

In long-run, full control would:

\begin{itemize}
    \item Reduce welfare
    \begin{itemize}
        \item Free market rent $>$ average controlled rent $\rightarrow$ reduce rental company's marginal rent
        \item Less quantity of rental housing supplied $\rightarrow$ higher initial and average rents in controlled market
    \end{itemize}
\end{itemize}

\vspace{0.5cm}

In long-run, no degree of rent control is optimal:

\begin{itemize}
    \item Less control $\rightarrow$ higher initial rents in controlled market to divert renters to free market
    \begin{itemize}
        \item Higher-income hhlds stay less in same unit $\rightarrow$ more sort to free market  
        \item Lower-income hhlds stay longer $\rightarrow$ pay more in controlled market
    \end{itemize}
\end{itemize}

\end{frame}
    
%%%%%%%%%%%%%%%%%%%%%%%%%
%%%%%%%%%%%%%%%%%%%%%%%%%
\begin{frame}{Road map}

\begin{itemize}
    \item History of rent control laws in Toronto
    \item Model 
    \item Calibration
    \item Experiments
    \item Conclusion
\end{itemize}

\end{frame}

%%%%%%%%%%%%%%%%%%%%%%%%%
%%%%%%%%%%%%%%%%%%%%%%%%%

{\setbeamercolor{background canvas}{bg=black}
\begin{frame}{}

\centering \Large
\textcolor{white}{History of rent control in Toronto}
    
\end{frame}

}

\begin{frame}{A timeline of rent control laws in Toronto}
  \centering
     \includegraphics[width=1\textwidth]{figures_for_slides/history_of_rent_control_Toronto.png}
\end{frame}

{\setbeamercolor{background canvas}{bg=black}
\begin{frame}{}

\centering \Large
\textcolor{white}{Model}
    
\end{frame}

}

%%%%%%%%%%%%%%%%%%%%%%%%%
%%%%%%%%%%%%%%%%%%%%%%%%%
\begin{frame}{Model}{Household preferences and demographic structure}
Household preferences:
\begin{equation*}
    u(c,h)=\frac{\left(c^{1-\gamma} h^{\gamma}\right)^{1-\sigma}}{1-\sigma},
\end{equation*}

\begin{itemize}
    \item $c$: non-housing consumption
    \item $h$: housing consumption
\end{itemize}

\hspace{0.5cm}

Overlapping generations of finitely-lived households:
\begin{itemize}
    \item Live from $j=1$ to $j=J$ and die with certainty after $j=J$
    \item Work until age $j<R$ and retire at age $j=R$
\end{itemize}

\end{frame}

\begin{frame}{Model}{Household heterogeneity}
Households are heterogeneous in:
\begin{itemize}
    \item Age, $j \in \mathcal{N}$
    \item Financial wealth, $a \in \mathcal{R_{+}}$
    \item Housing tenure, $o \in \{0,1\}$, $0=$renter, $1=$owner
    \item Rental market in which they rent, $f \in \{0,1\}$, $0=$controlled market, $1=$free market
    \item Years in same unit, $\iota \in \mathcal{N}$
    \item Housing size, $h \in \mathcal{R_{++}}$
    \item Mortgage debt, $m \in \mathcal{R_{+}}$
    \item Labor productivity, $e \in \mathcal{R_{++}}$
    \item $s$ to denote household state $(a,o,f,\iota,h,m,e)$
\end{itemize}

\end{frame}

\begin{frame}{Model}{Household incomes}

\begin{itemize}
    \item Labor income, $y_{\ell,j}(e)=W \zeta_{j} e$, for $j<R$ $\zeta_{j}$: deterministic age-dependent productivity
    \begin{itemize}
        \item Retirees receive pension, $y_{R}(e)$, $j\geq R \rightarrow$ higher productivities receive higher pensions
        \item Intra-generational persistence of $e:$, $\log e' = \rho_{e} \log e + \epsilon_{e},\quad \epsilon_{e}\sim N(0,\sigma_{e})$
        \item Inter-generational persistence of $e:$, $\log e^{\text{child}} = \hat{\rho}_{e} \log e^{\text{parent}} + \hat{\epsilon}_{e},\quad \hat{\epsilon}_{e}\sim N(0,\hat{\sigma}_{e})$
    \end{itemize}
    \item Capital income, $y_{k}(a)=ar$, $r$: risk free rate
    \item Taxable income, $y^{x}_{j}(e,a)=y_{\ell,j}(e)+\nu y_{k}(a)$, $\nu$: taxable portion of capital income
    \item Progressive taxes, $y^{\tau}_{j}(e,a) = y^{x}_{j}(e,a)-\tau_{\ell}y^{x}_{j}(e,a)^{\psi}$, $\tau_{\ell}$: avg. tax, $\psi$: progressivity
\end{itemize}
\end{frame}

%%%%%%%%%%%%%%%%%%%%%%%%%
%%%%%%%%%%%%%%%%%%%%%%%%%

\begin{frame}{Model}{Renters (part 1)}

\begin{itemize}
    \item Renters choose to rent in free market $f=1$, or controlled market $f=0$
    \begin{itemize}
        \item If $f=1$, pay $p^{r}_{1}h$, $p^{r}_{1}$: free market rent
        \item If $f=0$, pay $\frac{p^{r}_{0}}{(1+\kappa)^{\iota}} h$, $p^{r}_{0}$: controlled market rent of new tenants
        \begin{itemize}
            \item $\iota$: number of periods in the same unit
            \item $\kappa$: rate of decline in real controlled rent with each additional period tenant remains in same unit
        \end{itemize}
    \end{itemize}
\end{itemize}

\end{frame}

\begin{frame}{Model}{Renters (part 2)}
  \centering
    \begin{columns}[T]
    \begin{column}{0.6\textwidth}  
     \caption{Qualitative depiction of rent path based on rental market choice and years in same rental}
     \includegraphics[width=1.4\textwidth]{figures/path_of_rent_diagram.pdf}
     \end{column}
    \begin{column}{0.4\textwidth}
     \begin{itemize}
        \item To clear markets, renters in controlled market pay initially more than in free market
        \begin{itemize}
            \item High-income renters plan to buy early on $\rightarrow$ sort to free market to avoid premium
            \item Low-income renters plan to remain renters longer $\rightarrow$ pay premium and enjoy lower average rents in controlled market
        \end{itemize}
    \end{itemize}
    \end{column}
    \end{columns}
\end{frame}

\begin{frame}{Model}{Homeowners}

\begin{itemize}
    \item Owners have access to long-term mortgages, $m$
    \item At origination pay $\omega_{1}$ and must satisfy:
    \begin{itemize}
        \item LTV constraint $m \leq \lambda_{m,1}(ph) p h$
        \item GDS constraint $\frac{\tau_{p}ph+(r_{m}+\upsilon_{3})m}{y_{\ell,j}(e)+y_{k}(a)} \leq \lambda_{m,2}\left(\frac{m}{ph}\right)$
    \end{itemize}
    \item Pay at least $\upsilon_{1}m$ each period or incur cost of origination again
    \item Can pay up to $\upsilon_{2}>\upsilon_{1}$ without prepayment penalty
    \begin{itemize}
        \item Otherwise pay prepayment penalty: $\tau_{pp}=\omega_{2}\mbox{max}\left\{(1-\upsilon_{2})m-m',0\right\}$
    \end{itemize}
    \item Must fully settle debt when $m\leq \upsilon_{4}ph$
\end{itemize}

\end{frame}

\begin{frame}{Model}{Housing}

\begin{itemize}
    \item All housing depreciates at a rate $\delta_{h}$
    \item All housing is subject to a property tax rate $\tau_{p}$
    \item Housing receives exogenous capital gains $g_{p}$
    \begin{itemize}
        \item High priced cities $\rightarrow$ (rent$<$ongoing operation costs) $\rightarrow$ no supply of rental housing
        \item Capital gains tilt the user-cost equation to incentivize supply of rental housing
    \end{itemize}
\end{itemize}

\end{frame}

%%%%%%%%%%%%%%%%%%%%%%%%%
%%%%%%%%%%%%%%%%%%%%%%%%%
\begin{frame}{Model}{Dynamic program}
\footnotesize
\begin{align}
  \label{eq:vf}
  V_{j}(s)=\max_{c,a',o',f',\iota',h',m'}\ \biggl\{u\left(c,h\right)+
  \beta \mathbb{E}\biggl[V_{j+1}(s'|e)\biggl]\biggr\}
\end{align}
subject to
\begin{flalign}
  \begin{split} \label{eq:bc}
  & c+a'+(1-o')\left[(1-f')\frac{p^{r}_{0}}{(1+\kappa)^{\iota'-1}}h'+f'p^{r}_{1}h'\right]+o'(ph'-m')+o\left(ph\left(\tau_{p}+\delta_{h}\right)+r_{m}m\right) + \\ & \quad \tau_{b}ph'(\mbox{purchase house=1})+\tau_{s}ph(\mbox{sell house=1})+ \omega_{1}(\mbox{originate mortgage=1})+\\ & \quad \tau_{pp}(\mbox{prepay more than allowed=1}) = \mbox{(after tax income from all sources)}+a+o(ph-m)
  \end{split} \\
  & o' \in \{0,1\}, \hspace{0.2cm} f' \in \{0,1\} \hspace{0.2cm} \mbox{if} \hspace{0.2cm} o'=0 \label{eq:rental_market_type} \\
  & h' \geq \underline{h}_{r} \hspace{0.2cm} \mbox{if} \hspace{0.2cm} o'=0, \hspace{0.2cm}, \hspace{0.2cm} h' \geq \underline{h}_{o} \hspace{0.2cm} \mbox{if} \hspace{0.2cm} o'=1  \label{eq:min_housing}\\
  & \mbox{Satisfy LTV \& GDS constraints if originating new mortgage}
\end{flalign}

Policy functions:
\begin{itemize}
    \item $c_{j}(s)$, $a'_{j}(s)$, $o'_{j}(s)$, $f'_{j}(s)$, $\iota'_{j}(s)$, $h'_{j}(s)$, and $m'_{j}(s)$
\end{itemize}

\end{frame}

%%%%%%%%%%%%%%%%%%%%%%%%%
%%%%%%%%%%%%%%%%%%%%%%%%%
\begin{frame}{Model}{Rental company (part 1)}
Rental company buys housing and rents them out: 
\begin{align}\label{eq:rental_company_problem}
V(X)=\underset{X'_{0},X'_{1}} {\text{$\mbox{max} $}}\left\{\sum_{i=0}^{i} \overline{p}^{r}_{i}X'_{i}-c_{m}X'^{\theta}-p X'+p(1-\delta_{h}-\tau_{p}+g_{p})X+\frac{1}{1+r}V(X')\right\}.
\end{align}
subject to
\begin{align*}
& X'_{0}=\mbox{min}\{\underline{X}_{0},X'\}, \hspace{0.2cm} X'_{1} \geq 0, \hspace{0.2cm} X'=\sum^{1}_{i=0} X'_{i} 
\end{align*}

\begin{itemize}
    \item $\underline{X'}_{0}$: amount of mandated controlled units; $X'_{1}$: supply on free market
    \item $\overline{p}^{r}_{1}=p^{r}_{1}$: free market rent, which is also the marginal rent for rental company
    \item $\overline{p}^{r}_{0}$: average rent in controlled market
    \item $c_{m}$: avg. management cost; $\theta$: convexity of cost of supplying more rental housing
\end{itemize}

\end{frame}

\begin{frame}{Model}{Rental company (part 2)}

\begin{itemize}
    \item Inverse rental supply function is therefore: $\overline{p}^{r}_{1}=c_{m}\theta X'^{\theta-1}+p-\left\{\frac{1-\delta_{h}-\tau_p+g_{p}}{1+r}\right\}p$ 
    \begin{itemize}
        \item When rental markets are fully controlled, avg rents in controlled market $\overline{p}^{r}_{0}$ replace $\overline{p}^{r}_{1}=p^{r}_{1}$
        \item In eq'm $\overline{p}^{r}_{1}>\overline{p}^{r}_{0}$ $\rightarrow$ moving to full control discontinuously reduces marginal rent and supply
    \end{itemize}
\end{itemize}

\end{frame}

%%%%%%%%%%%%%%%%%%%%%%%%%
%%%%%%%%%%%%%%%%%%%%%%%%%
\begin{frame}{Model}{Housing construction company}

\begin{itemize}
    \item Competitive company constructs new housing with technology, $TC(I)=c_{1}\left(\frac{1}{\frac{1}{\epsilon}+1}\right)I^{\frac{1}{\epsilon}+1}$
    \item Optimal investment: $I=\left(\frac{p}{c_{1}}\right)^{\epsilon}$; $p$: housing price; $\epsilon$: rental supply elasticity; $c_{1}$: scalar
    \item Law of motion for housing $H'=(1-\delta_{h})H+I$
\end{itemize}

\end{frame}

\begin{frame}{Model}{Government}

Must balance budget each period:
\begin{itemize}
    \item $G + \sum_{j=R}^J \int_\mathcal{S} y_{R}(s) \mathrm{d}\Psi(s) + \sum_{j=1}^J \int_\mathcal{S} T \mathrm{d}\Psi(s) = \sum_{j=1}^J \int_\mathcal{S}y^{\tau}_{j}(s) \mathrm{d}\Psi(s)+\tau_p p H$
    \begin{itemize}
        \item $G:$ public good; picked in benchmark to clear budget, held fixed in counterfactuals
        \item $T:$ lump sum transfer; set to 0 in benchmark, picked in counterfactuals to clear budget
    \end{itemize}
\end{itemize}

\end{frame}

\begin{frame}{Model}{Equilibrium definition}

A stationary eq'm in the model is:

\begin{itemize}
    \item a sequence of aggregate prices and quantities, $\left\{H',I,\{X'_{i},\overline{p}^{r}_{i}\}^{1}_{i=0},p\right\}$;
    \item value and policy functions, $\biggl\{V_{j}(\cdot),c_{j}(\cdot),a'_{j}(\cdot),o'_{j}(\cdot),f'_{j}(\cdot)\iota'_{j}(\cdot),h'_{j}(\cdot),m'_{j}(\cdot)\biggl\}_{j=1}^J$; and
    \item a sequence of measures, $\left\{\Psi_{j}(\cdot)\right\}_{j=1}^J$, that:
\end{itemize}
\begin{enumerate}
%\item all nominal prices rise at the rate of $i$, but real prices remain fixed;
\item Solve the household's dynamic problem subject to constraints;
\item Solve the rental company's problem subject to constraints;
\item Satisfy the construction company's first-order condition;
\item Satisfy the government's budget constraint;
\item Clear the housing market;
\item Clear the rental housing markets;
\item And satisfy the laws of motion for the distribution of households
\end{enumerate}
    
\end{frame}

{\setbeamercolor{background canvas}{bg=black}
\begin{frame}{}

\centering \Large
\textcolor{white}{Calibration}
    
\end{frame}

}

%%%%%%%%%%%%%%%%%%%%%%%%%
%%%%%%%%%%%%%%%%%%%%%%%%%
\begin{frame}{Calibration}{External calibration (part 1)}

\input{tables_for_slides/calibration_1.tex}

\end{frame}

\begin{frame}{Calibration}{External calibration (part 2)}

\input{tables_for_slides/calibration_2.tex}

\end{frame}

\begin{frame}{Calibration}{External calibration (part 3)}

\input{tables_for_slides/calibration_3.tex}

\end{frame}

\begin{frame}{Calibration}{Internal calibration}

\input{tables_for_slides/calibration_4.tex}

\end{frame}

\begin{frame}{Calibration}{Non-targeted moments - aggregates (part 1)}

\input{tables_for_slides/aggregate_non_targeted_statistics_1.tex}

\end{frame}

\begin{frame}{Calibration}{Non-targeted moments - aggregates (part 2)}

\input{tables_for_slides/aggregate_non_targeted_statistics_2.tex}

\end{frame}

\begin{frame}{Calibration}{Non-targeted moments - lifecycle}
  \centering
     \includegraphics[width=0.9\textwidth]{figures/lifecycle_plots.pdf}
\end{frame}

\begin{frame}{Calibration}{Non-targeted moments - tenure by productivity}
  \centering     \input{tables_for_slides/tenure_by_productivity_benchmark_1.tex}
\end{frame}

%%%%%%%%%%%%%%%%%%%%%%%%%
%%%%%%%%%%%%%%%%%%%%%%%%%

{\setbeamercolor{background canvas}{bg=black}
\begin{frame}{}

\centering \Large
\textcolor{white}{Experiments}
    
\end{frame}

}

\begin{frame}{Experiments}{Overview}

We conduct 4 experiments that we compare to the benchmark:
\begin{itemize}
    \item Long-run full rent decontrol
    \item Long-run full rent control
    \item Optimal degree of rent control in long-run
    \begin{itemize}
        \item Maximize ex-ante expected utility of newborns \citep{conesa-et-al,gkkoc}
    \end{itemize}
    \item Transition from benchmark to long-run full rent decontrol - permanent "MIT shock"
\end{itemize}

\vspace{0.5cm}

In all our experiments, we:
\begin{itemize}
    \item Keep tax system fixed
    \item Balance government budget with lump sum transfers
\end{itemize}
    
\end{frame}

\begin{frame}{Experiments}{Measure of rents}

\begin{enumerate}
    \item \underline{Free market rent}: rent in free market.
    \item \underline{Average controlled rent}: avg. rent in controlled market.
    \item \underline{Rent of new tenants}: rent new tenants pay when moving into a rental unit
    \item \underline{Average rent}: weighted avg. of rents in controlled and free markets
    \item \underline{Marginal rent}: rent that rental company receives if it supplied an additional rental unit
    \begin{itemize}
        \item When free markets operate, marginal rent is free market rent
        \item When all rental units are controlled, marginal rent is avg. rent in controlled market
    \end{itemize}
\end{enumerate}

\end{frame}

\begin{frame}{Experiments}{Experiment 1: Long-run full decontrol, aggregate outcomes}

\begin{columns}[T]
\begin{column}{0.5\textwidth}
\input{tables_for_slides/aggregate_effects_1.tex}
\end{column}

\begin{column}{0.53\textwidth}
    \begin{itemize}
    \item Full decontrol $\rightarrow$ reduced incentive to remain in same unit $\rightarrow$ more homeownership
    \item Less rental demand $\rightarrow$ lower free market rent
    \item Lower rents of new tenants $\rightarrow$ benefits low-income, young hhlds with highest MU $\rightarrow$ welfare gain of newborns
    \item In benchmark free market rent is 8.8\%$>$avg controlled rent $\rightarrow$ higher avg. rent under full decontrol
    \item Higher avg. rent $\rightarrow$ hurts renters already under rent control $\rightarrow$ welfare loss when considering everyone
    \end{itemize}
\end{column}

\end{columns}

\end{frame}

\begin{frame}

\begin{figure}[H]
  \centering
    \caption{Qualitative depiction of impact of full rent decontrol}
    \label{fig:diagram_full_decontrol}
    \vspace{-0.26cm}
    \includegraphics[width=0.95\textwidth]{figures/full_rent_decontrol_diagram.pdf}
\end{figure}

\end{frame}

\begin{frame}{Experiments}{Experiment 1: Long-run full decontrol, distributional outcomes}

\input{tables_for_slides/distributional_impacts_1.tex}

\begin{itemize}
    \item $e_{1}$ are lowest income $\rightarrow$ value lower initial rents more than lower average rents
    \item $e_{2}-e_{3}$ remain long enough to benefit from rent control in same units and have higher incomes $\rightarrow$ value lower avg. rents more than lower initial rents
    \item $e_{4}-e_{7}$ do not remain long enough to benefit from rent control $\rightarrow$ value lower initial rents 
    \item Rent effect is first order, lump sum transfer second order, house price third order
\end{itemize}

\end{frame}

\begin{frame}{Experiments}{Experiment 2: Long-run full control, aggregate outcomes}

\begin{columns}[T]
\begin{column}{0.5\textwidth}
\input{tables_for_slides/aggregate_effects_2.tex}
\end{column}
\begin{column}{0.53\textwidth}
    \begin{itemize}
    \item Full control $\rightarrow$ incentive for more hhlds to remain longer in same rental 
    \item Lower marginal rents $\rightarrow$ reduce rentals supplied $\rightarrow$ increase avg. rents    
    \item Newborns face higher rents $\rightarrow$ reduce welfare 
    \item Why does homeownership increase also under full control?
    \begin{itemize}
        \item Under full control, lower house prices cause homeownership to rise
        \item Under full decontrol, homeownership causes house prices to rise
    \end{itemize}
    \end{itemize}
\end{column}

\end{columns}
\end{frame}

\begin{frame}

\begin{figure}[H]
  \centering
    \caption{Qualitative depiction of impact of full rent control}
    \label{fig:diagram_full_control}
    \vspace{-0.26cm}
    \includegraphics[width=0.97\textwidth]{figures/full_rent_control_diagram.pdf}
\end{figure}

\end{frame}

\begin{frame}{Experiments}{Experiment 3: Optimal degree of rent control in long-run}

Is little rent control better? We find the answer is "No" and that rental market sorting is key 
\begin{itemize}
    \item Less control $\rightarrow$ higher initial rents in controlled market to divert renters to free market
    \begin{itemize}
        \item Higher-income hhlds stay less in same unit $\rightarrow$ more sort to free market  
        \item Lower-income hhlds with high MU stay longer $\rightarrow$ pay more in controlled market
    \end{itemize}
\end{itemize}

\end{frame}

\begin{frame}{Experiments}{Experiment 4: Transition from benchmark to full rent decontrol}
  \centering
     \includegraphics[width=0.9\textwidth]{figures/transition_plots.pdf}
\begin{itemize}
    \item Decontrolling $\rightarrow$ raises rents for renters under control for a while $\rightarrow$ less rental demand 
    \item Previously-uncontrolled rents decline upon decontrolling $\rightarrow$ increase newborns' welfare 
    \item Less incentive to rent $\rightarrow$ accumulate more wealth to buy $\rightarrow$ house prices rise
    \item Higher house prices $\rightarrow$ more property tax revenue $\rightarrow$ larger lump sum 
    \item More wealth $\rightarrow$ more capital income $\rightarrow$ existing renters demand larger rentals $\rightarrow$ rents bounce back a bit
\end{itemize}     
     
\end{frame}

{\setbeamercolor{background canvas}{bg=black}
\begin{frame}{}

\centering \Large
\textcolor{white}{Conclusion}
    
\end{frame}

}

%%%%%%%%%%%%%%%%%%%%%%%%%
%%%%%%%%%%%%%%%%%%%%%%%%%
\begin{frame}{Conclusion}{Part 1}

Build GE model with controlled and free rental markets to examine impact of rent control on:
\begin{itemize}
    \item Aggregate outcomes
    \item Distributional outcomes
    \item Short-run outcomes
\end{itemize}

\vspace{0.5cm}

We calibrate the model to data from Toronto:
\begin{itemize}
    \item Rental markets are partly controlled and partly free
    \item Model does reasonably well at imitating aggregate, life-cycle, and distributional data
    \item Use calibrated model to conduct a series of experiments
\end{itemize}

\end{frame}

\begin{frame}{Conclusion}{Part 2}

We find that:
\begin{itemize}
    \item In long-run, full rent decontrol would:
    \begin{itemize}
        \item Increase avg. rents $\rightarrow$ hurt renters who already live in controlled units
        \item Reduce rents of new tenants $\rightarrow$ benefit high- and low-income, young households
    \end{itemize}
    \item In short-run, full rent decontrol would:
    \begin{itemize}
        \item Immediately reduce free market rents and initial rents of new tenants $\rightarrow$ welfare gains
    \end{itemize}
    \item In long-run, full rent control would:
    \begin{itemize}
        \item Lower rental company's marginal rent $\rightarrow$ less rentals supplied $\rightarrow$ higher initial and avg. rents
    \end{itemize} 
    \item No degree of rent control is optimal in long-run:
    \begin{itemize}
        \item Less control $\rightarrow$ higher initial rents in controlled market to divert renters to free market
        \item Higher-income hhlds stay less in same unit $\rightarrow$ more sort to free market  
        \item Lower-income hhlds stay longer $\rightarrow$ pay more in controlled market
    \end{itemize}    
\end{itemize}
    
\end{frame}

\begin{frame}{Conclusion}{Limitations and future research}

Due to computational limitations, we abstract from:
\begin{itemize}
    \item Aggregate uncertainty
    \begin{itemize}
        \item Important for studying benefits of rent rigidity with rent control during RBC fluctuations
    \end{itemize}
    \item Endogenous labor supply
    \begin{itemize}
        \item Key to exploring rent control with subsidies to low-income, young hhlds, paid by income taxes
    \end{itemize}
\end{itemize}
 
\end{frame}

{\setbeamercolor{background canvas}{bg=black}
\begin{frame}{}

\centering \Large
\textcolor{white}{References}
    
\end{frame}

}

\bibliographystyle{aer}
\bibliography{reference_files/bib}

\end{document}
